\documentclass[11pt]{article}
\usepackage[margin=1in]{geometry}
\usepackage{amsmath,amssymb,booktabs,microtype}
\usepackage{xcolor}
\usepackage[colorlinks=true,linkcolor=blue,urlcolor=blue]{hyperref}
\newcommand{\Sym}{\operatorname{Sym}}
\newcommand{\Gr}{\operatorname{Gr}}
\newcommand{\cS}{\mathcal S}
\title{Grassmann Schemes: Permutation Molien Formula\\Proof and Exact Verification}
\author{Executable companion note}
\date{17 July 2026}
\begin{document}
\maketitle
\begin{abstract}
For $GL_N(q)$ acting on the $k$-subspaces of $\mathbb F_q^N$, we prove the
permutation Molien formula and verify its invariant-subspace input by direct
prime-field matrix construction.  Three representative cases test points,
hyperplanes, and a nontrivial scalar kernel.
\end{abstract}

\section{General formula}
Let $X=\Gr_q(N,k)$ and $V=\mathbb C[X]$.  If $c_\ell(g)$ is the number of
$\ell$-cycles of $g$ on $X$, the cycle-block determinant identity gives
\begin{equation}\label{eq:master}
\cS_{G,X}(\mathbf t)=\frac1{|G|}\sum_{g\in G}
\prod_{i,\ell}(1-t_i^\ell)^{-c_\ell(g)}.
\end{equation}
The monomial basis of $\Sym^d(V)$ is indexed by $d$-multisets of subspaces,
so Burnside's lemma also proves
\begin{equation}\label{eq:orbit}
[m_\tau]\cS_{G,X}=\#\left(G\backslash\prod_j\Sym^{\tau_j}(X)\right).
\end{equation}

For a linear matrix $g$, define
\begin{equation}\label{eq:fixed}
F_r^{\rm Gr}(g)=\#\{U\leq\mathbb F_q^N:\dim U=k,\ g^rU=U\}.
\end{equation}
Every $\ell$-cycle contributes its $\ell$ vertices to $F_r$ exactly when
$\ell\mid r$.  Möbius inversion therefore yields
\begin{equation}\label{eq:mobius}
c_\ell^{\rm Gr}(g)=\frac1\ell\sum_{d\mid\ell}
\mu(\ell/d)F_d^{\rm Gr}(g).
\end{equation}
Together, \eqref{eq:master}, \eqref{eq:fixed}, and \eqref{eq:mobius} prove the
claimed exact formula.  Interpreting the ambient space as an
$\mathbb F_q[x]$-module with $x$ acting as $g^r$ identifies
\eqref{eq:fixed} with the stated invariant-submodule problem.

Pair orbits are classified by $\dim(U\cap W)$, and so, with
$s=\min(k,N-k)$,
\[
[m_{(1)}]=1,\qquad [m_{(2)}]=[m_{(1,1)}]=s+1.
\]

\section{Verification strategy}
All matrices in $GL_N(q)$ are enumerated over the prime field.  Every
$k$-subspace is generated and stored in unique reduced-row-echelon form.
Matrix multiplication then gives both the induced vertex permutation and the
independent invariance test in \eqref{eq:fixed}.  Scalar matrices are deduplicated
before averaging, so the average is over the actual permutation image; this is
equivalent to averaging over $GL_N(q)$ because the scalar kernel has constant
size.

The target coefficient is computed from actual permutation matrices using
Newton's trace recurrence for $h_d(P)=\operatorname{tr}(\Sym^dP)$.  The same
coefficient is independently extracted from \eqref{eq:master} and counted as
the explicit orbit number in \eqref{eq:orbit}.  Finally, invariant-subspace
counts for powers of every matrix representative must reconstruct every cycle
through \eqref{eq:mobius}.

\section{Exact results}
The coefficient vector is ordered by
$\tau=(1),(2),(1,1),(3),(2,1),(1,1,1)$.
\begin{center}
\begin{tabular}{@{}lrrrl@{}}
\toprule
Action & $|GL|$ & $|G\text{ image}|$ & $|X|$ & common vector\\
\midrule
$\Gr_2(3,1)$ & 168 & 168 & 7 & $(1,2,2,4,5,6)$\\
$\Gr_2(3,2)$ & 168 & 168 & 7 & $(1,2,2,4,5,6)$\\
$\Gr_3(2,1)$ & 48 & 24 & 4 & $(1,2,2,3,4,5)$\\
\bottomrule
\end{tabular}
\end{center}
All 18 exact triple comparisons and all invariant-subspace reconstructions
pass.  The final row explicitly checks the scalar kernel.  In the first case,
the direct determinant and cycle averages at $(0.037,0.061)$ both equal
$1.1152638116688691$ (reported error $0$).

\section{Reproduction}
Run
\begin{verbatim}
marimo run marimo_notebooks/grassmann_scheme.py
\end{verbatim}
from the repository root.  The implementation is intentionally restricted to
prime fields for transparency; the theorem itself is not so restricted.
\end{document}
